\section{Introduction}\label{sec:intro}

Diffusion transformers trained at about one megapixel are now asked for $4096\times4096$ images, which hold \N{c@tokratio} times more tokens than any training image. Training-free methods handle this by editing the rotary position embedding (RoPE)~\citep{su2024roformer} at inference, without new weights or extra sampling passes. They rescale its frequencies~\citep{chen2023pi,bloc97ntk,peng2024yarn}, schedule the rescaling with the noise level~\citep{issachar2025dype} or sharpen attention with a global logit temperature~\citep{jin2023training,nakanishi2025ssmax}. The most recent methods also scale each rotary band by a gain read from the spectrum of the current latent~\citep{sega2026}. Each is compared as a package with its predecessor, and the gain is credited to its newest component. This paper asks which edit produces the gain, and whether a per-band law with a principled basis can do better.

\begin{wraptable}{r}{0.45\linewidth}
\vspace{-12pt}
\caption{\textbf{Knobs each arm turns.} The shaded row is the reference, \segaofficial; gray marks a context baseline. The factorial of \cref{sec:decomp} adds \ntk{} with the temperature and the \sega{} base alone. \Cref{tab:methods} (\cref{app:impl}) expands the table with formulas and sample sizes. Key result: only the per-band column separates \tempsb{} from \segaofficial.}
\label{tab:knobs}
\centering\scriptsize
\setlength{\tabcolsep}{2.4pt}
\begin{tabular}{@{}l l c l@{}}
\toprule
Arm & Base & Temp. & Per-band \\
\midrule
\ntk & \ntk & \no & \no \\
\textcolor{ctxgray}{\dype} & \textcolor{ctxgray}{ramps in $\sigma$} & \yes & \no \\
\rowcolor{papertint}\segaofficial & \sega & \yes & \yes{} sensor \\
\tempsb & \sega & \yes & \no \\
\constancytm & \sega & \yes & \yes{} gated \\
\segactm & \sega & \yes & \yes{} both \\
\bottomrule
\end{tabular}
\vspace{-8pt}
\end{wraptable}
Package comparisons cannot answer that question, because each method turns several of the knobs in \cref{tab:knobs} at once. \sega{} raises the frequency base, adds a global temperature and adds a content-adaptive per-band allocation in one step, and its paper credits the improvement to the allocation. A comparison against \dype{} or an NTK base cannot tell which of the three did the work. The same confound affects any new per-band law that is tested only against complete methods. Evaluation adds a second difficulty: the reward models that rank these methods see a 224-pixel view, which cannot resolve the finest octaves of a 4K image.

We introduce a controlled decomposition of \sega~\citep{sega2026}, a strong recent method with released code, on the checkpoint its authors use (\krea). Its central control, \tempsb, keeps the released frequency base, applies \sega's reference logit gain $s^{0.16}$ to every query and key dimension ($s=\N{c@s}$ per side) and drops only the allocation. A two-by-two factorial then crosses the base (\ntk{} or \sega) with the temperature (none or $s^{0.16}$). Base, temperature and allocation thus add up to the gain of \sega{} over an \ntk{} base. All arms share prompts, seeds and sampler and differ only in the frequencies and gains of the attention logit. We fixed a dev/held-out split before any held-out score was seen and registered every follow-up test before generating its images. Under the same protocol we test four per-band controllers derived from efficient coding and contrast constancy. We do not propose a new state-of-the-art method.

We find that the global temperature does the work (\cref{fig:hero}). On \N{st@f@n} held-out prompts it raises ImageReward by \N{st@f@a@temp@ir} \N{st@f@ci@temp@ir}, the whole \N{st@f@tot@ir} gain of \sega{} over \ntk. The \sega{} base adds \N{st@f@d@base@ir} \N{st@f@ci@base@ir}, and no base or interaction test survives Holm correction on any of ten metrics. The allocation adds \N{st@h@d@musiq} \N{st@h@ci@musiq} full-frame MUSIQ on \N{st@e@n} fresh prompts and no Holm-surviving reward-model gain. On \N{st@h@n@ir} pooled prompts, \tempsb{} is equivalent to released \sega{} within preregistered margins on ImageReward (\N{st@h@d@ir} \N{st@h@ci@ir}) and PickScore, but the pool reuses the \N{st@f@n} held-out prompts we had already inspected. On the fresh prompts alone, PickScore differs by \N{st@hf@d@pick} \N{st@hf@ci@pick}, crossing its $\pm\N{st@h@m@pick}$ margin. No multiplier in a fixed temperature sweep beats $s^{0.16}$ on a reward model or full-frame MUSIQ, although 1.45 wins crop MUSIQ under Holm (\N{st@s@d@t145@musiqc}). The principled laws tie or lose to released \sega{} on every reward and image-quality metric; a tie means that the paired 95\% interval covers zero. On eight dev prompts at matched guidance, \fluxdev{} carries \N{st@k@two@ckg35} times \N{st@k@twoci@ckg35} the lattice power of \krea, so an earlier lattice failure is a checkpoint effect. On the same prompts, the diffusers default flow shift changes ImageReward by \N{st@ds@d@def@ir} \N{st@ds@ci@def@ir}.

To summarize, we decompose a strong training-free 4K method under a registered protocol and find that its global attention temperature, rather than its base or its allocation, produces its reward-model gain. We highlight our main contributions below.
\begin{itemize}
\item We introduce a factorial decomposition and a paired, registered protocol that attribute the gains of training-free RoPE methods at 4K to individual edits.
\item We derive four per-band controllers from efficient coding and contrast constancy and find that the two evaluated held out tie or lose to released \sega.
\item We show that the global temperature, not the base, carries \sega's whole ImageReward gain over an \ntk{} base, \N{st@f@a@temp@ir} \N{st@f@ci@temp@ir}, and that the allocation's full-frame MUSIQ gain replicates on fresh prompts.
\end{itemize}
